% ============================================================
%  MODULE I — Introduction to Financial Markets
% ============================================================

\chapter{Introduction to Financial Markets}

\begin{notebox}
This module builds the conceptual foundation. Every advanced topic later — options pricing,
order flow, macro positioning — assumes fluency with what is covered here.
\end{notebox}

% ────────────────────────────────────────────────────────────
\section{What Are Financial Markets?}

A \textbf{financial market} is any organised or informal system in which buyers and sellers
exchange financial instruments — assets that represent a claim on future cash flows, ownership,
or contractual rights.

Markets exist to solve two universal problems:
\begin{enumerate}[label=\textbf{\arabic*.}]
  \item \textbf{Capital allocation} — directing money from those who have it (savers, investors)
        to those who need it (companies, governments, individuals).
  \item \textbf{Price discovery} — aggregating dispersed, private information from millions of
        participants into a single observable price.
\end{enumerate}

\begin{keyterm}{Financial Instrument}
Any contract that creates a financial asset for one party and a financial liability (or equity)
for another. Examples: shares, bonds, futures contracts, options, currencies, loans.
\end{keyterm}

% ────────────────────────────────────────────────────────────
\section{The Market Hierarchy}

\begin{center}
\begin{tabular}{>{\bfseries}p{4cm} p{9cm}}
\toprule
Market & What Trades Here \\
\midrule
Equity Markets & Ownership stakes (shares/stocks) in companies \\
Debt / Bond Markets & Loans to governments and corporations \\
Money Markets & Very short-term debt (\textless 1 year); cash equivalents \\
Foreign Exchange (FX) & Currency pairs \\
Commodity Markets & Physical goods: oil, gold, wheat, etc. \\
Derivatives Markets & Contracts derived from underlying assets (futures, options, swaps) \\
Crypto Markets & Digital/cryptographic assets \\
\bottomrule
\end{tabular}
\end{center}

\begin{notebox}
These markets are deeply interconnected. A rising dollar (\$) usually pressures commodity prices
(they are dollar-denominated). Rising interest rates affect stock valuations via the discount rate.
You will not understand equities fully without at minimum a passing knowledge of rates and FX.
\end{notebox}

% ────────────────────────────────────────────────────────────
\section{Primary vs. Secondary Markets}

\begin{keyterm}{Primary Market}
Where new securities are \emph{created and sold for the first time}.
A company doing an IPO (Initial Public Offering) is selling shares in the primary market.
Proceeds go \textbf{to the issuer}.
\end{keyterm}

\begin{keyterm}{Secondary Market}
Where previously issued securities are \emph{traded between investors}.
When you buy Apple stock on the NYSE, you are buying from another investor in the secondary market.
Proceeds go \textbf{to the seller}, not to Apple.
\end{keyterm}

This distinction matters: most people operate entirely in secondary markets, yet it is primary
market activity (IPOs, bond issuance) that directly funds the real economy.

% ────────────────────────────────────────────────────────────
\section{Why Prices Move}

Price is determined by \textbf{supply and demand}. In financial markets, that comes down to:

\begin{itemize}
  \item \textbf{Expectations} — prices are forward-looking. A stock represents the present value
        of all future cash flows, not current earnings alone.
  \item \textbf{New information} — earnings beats, central bank decisions, geopolitical shocks
        cause participants to re-assess expected future cash flows.
  \item \textbf{Sentiment and positioning} — markets overshoot fair value regularly because humans
        are emotional. Fear and greed are real forces.
  \item \textbf{Liquidity} — thin markets move more violently on the same information than deep ones.
\end{itemize}

\begin{keyterm}{Efficient Market Hypothesis (EMH)}
The theory that asset prices \emph{fully reflect all available information} at any given moment,
making it impossible to consistently outperform the market through skill alone.
There are three forms: \textbf{weak} (prices reflect past prices), \textbf{semi-strong}
(prices reflect all public information), and \textbf{strong} (prices reflect all information
including insider knowledge). Widely debated — practical traders generally reject strong form.
\end{keyterm}

% ────────────────────────────────────────────────────────────
\section{Return and Risk — The Foundational Trade-off}

\begin{formulabox}
\[
  \text{Return} = \frac{P_1 - P_0 + D}{P_0}
\]
where $P_0$ = entry price, $P_1$ = exit price, $D$ = dividends (or other income received).
\end{formulabox}

Every additional unit of expected return requires accepting additional risk. This is the
\textbf{risk-return trade-off} — the single most important idea in all of finance.

Common risk types:
\begin{itemize}
  \item \textbf{Market risk} (systematic) — risk from broad market moves; cannot be diversified away.
  \item \textbf{Idiosyncratic risk} (unsystematic) — company-specific risk; can be reduced by diversification.
  \item \textbf{Liquidity risk} — inability to exit a position at the desired price.
  \item \textbf{Credit risk} — the counterparty may default.
  \item \textbf{Inflation risk} — real returns eroded by price level increases.
  \item \textbf{Tail risk} — low-probability, very high-impact events (black swans).
\end{itemize}

% ────────────────────────────────────────────────────────────
\section{Time Value of Money}

This is the mathematical backbone of valuation.

\begin{formulabox}
\textbf{Future Value:}
\[
  FV = PV \times (1 + r)^n
\]
\textbf{Present Value:}
\[
  PV = \frac{FV}{(1 + r)^n}
\]
where $r$ = discount rate (required return), $n$ = number of periods.
\end{formulabox}

\begin{examplebox}
If you expect to receive \$1{,}000 in 3 years and your required return is 8\%:
\[
  PV = \frac{1000}{(1.08)^3} = \frac{1000}{1.2597} \approx \$793.83
\]
You should be willing to pay at most \$793.83 today for that future \$1{,}000.
\end{examplebox}

% ────────────────────────────────────────────────────────────
\section{Key Vocabulary to Memorise}

\begin{center}
\begin{tabular}{>{\bfseries}p{3.5cm} p{9.5cm}}
\toprule
Term & Definition \\
\midrule
Bull market & Sustained price increases, typically $>$20\% from trough \\
Bear market & Sustained price declines, typically $>$20\% from peak \\
Correction & A 10--20\% pullback from a recent high \\
Rally & A short-term price increase within any trend \\
Volatility & The degree of price fluctuation, often measured as standard deviation \\
Liquidity & How easily an asset can be bought or sold without moving the price \\
Spread & Difference between bid (buy) price and ask (sell) price \\
Long & Owning an asset; profit if price rises \\
Short & Borrowed and sold an asset; profit if price falls \\
Leverage & Using borrowed money to amplify position size \\
Margin & Collateral deposited to support a leveraged position \\
Mark-to-market & Daily revaluation of positions at current market prices \\
\bottomrule
\end{tabular}
\end{center}

\warningbox{
Leverage amplifies \emph{both} gains and losses proportionally. A 2$\times$ leveraged position
loses 20\% when the underlying moves 10\% against you. Many beginners blow up accounts by
misunderstanding leverage.
}
